\fancyhead{}
\fancyfoot{}
\cfoot{\thepage}

\lhead{Discusión}

\chapter{Discusión}

El desarrollo del Sistema de Asistencia Basado en Reconocimiento Facial para la FPUNE permitió abordar una problemática frecuente en el ámbito académico: la gestión eficiente y confiable de la asistencia estudiantil.

Los resultados obtenidos evidencian que, con una configuración adecuada de parámetros como el umbral de similitud y el modelo de reconocimiento, es posible alcanzar un desempeño satisfactorio en precisión y velocidad. Modelos como YOLOv8n-Face y ArcFace demostraron un equilibrio óptimo entre exactitud y rendimiento, resultando viables para entornos con recursos limitados.

Entre las principales limitaciones se identificaron la sensibilidad a la iluminación y las expresiones faciales, además de la necesidad de disponer de múltiples imágenes por usuario. Estos aspectos abren líneas de mejora orientadas a incrementar la robustez del sistema y ampliar la base de datos.

En síntesis, el prototipo demostró ser una solución técnica factible y beneficiosa, capaz de reducir errores humanos y optimizar los procesos administrativos relacionados con el control de asistencia.



\section{Logros alcanzados con relación a los objetivos específicos}

\textbf{1. Seleccionar los algoritmos de reconocimiento facial}

El objetivo se cumplió satisfactoriamente mediante la selección de modelos adecuados para la detección y reconocimiento facial. Se optó por el uso de \textbf{YOLOv8n-Face} para la detección de rostros y \textbf{ArcFace} para la verificación de identidad, debido a su equilibrio entre exactitud, velocidad de procesamiento y tamaño del modelo.  
Los resultados experimentales presentados en el Capítulo~5 mostraron que esta combinación alcanzó una precisión promedio del 85.2\,\% y un tiempo de detección de 0.32\,s, evidenciando un desempeño estable y eficiente en condiciones reales de aula. Estas métricas confirman la idoneidad de los algoritmos seleccionados para lograr un reconocimiento facial confiable dentro del contexto académico.

\textbf{Conclusión:} El objetivo se alcanzó plenamente, estableciendo una base sólida para el desarrollo del sistema de asistencia.

\textbf{2. Seleccionar las herramientas tecnológicas para el desarrollo del sistema}

Se cumplió el objetivo al definir una infraestructura tecnológica compatible con los requerimientos del proyecto. Se utilizaron \textbf{Python}, \textbf{OpenCV} y \textbf{Flask} en el backend, junto con \textbf{PostgreSQL} para la base de datos y \textbf{React} en el frontend, logrando una arquitectura modular, escalable y de fácil mantenimiento.

\textbf{Conclusión:} El conjunto de herramientas seleccionadas permitió integrar eficazmente las etapas de captura, reconocimiento y almacenamiento de datos. \par

\textbf{3. Identificar los requisitos para el desarrollo del sistema de registro de asistencia para la FPUNE}

El objetivo se alcanzó mediante la recopilación de información sobre los procedimientos actuales de control de asistencia y las necesidades específicas de los docentes y administrativos. Se definieron los requisitos funcionales y no funcionales, incluyendo seguridad de datos, facilidad de uso, precisión del reconocimiento y disponibilidad en tiempo real.

\textbf{Conclusión:} La identificación de requisitos permitió diseñar un sistema ajustado a las condiciones técnicas y operativas de la FPUNE. \par

\textbf{4. Desarrollar el sistema de reconocimiento facial y gestión de asistencias}

Se cumplió el objetivo mediante el desarrollo del \textbf{Sistema de Asistencia Basado en Reconocimiento Facial para la FPUNE}, que integra los módulos de detección, reconocimiento y registro automático en la base de datos. El sistema fue implementado con una interfaz intuitiva, capaz de registrar la asistencia sin intervención manual.

\textbf{Conclusión:} El desarrollo del sistema demostró la viabilidad técnica del proyecto y su potencial de aplicación en entornos académicos reales. \par

\textbf{5. Realizar pruebas del funcionamiento del prototipo}

El objetivo se cumplió mediante la ejecución de pruebas experimentales en distintas condiciones de iluminación, distancia y ángulo de cámara, con el propósito de evaluar el desempeño real del sistema de reconocimiento facial. Los resultados obtenidos, presentados en el Capítulo~5, evidenciaron una precisión promedio del 85.2\,\%, un tiempo de detección de 0.32\,s por rostro y una tasa de registros automáticos correctos del 97.5\,\%. Estos valores cuantitativos reflejan un comportamiento estable del modelo ante variaciones comunes en el entorno de aula, confirmando su capacidad para operar de manera confiable en escenarios reales de uso.

\textbf{Conclusión:} Las pruebas demostraron que el sistema mantiene un rendimiento consistente y predecible dentro de los márgenes normales de iluminación y posición, validando su fiabilidad técnica para el registro automatizado de asistencias.

\textbf{6. Evaluar los resultados obtenidos}

El objetivo se desarrolló mediante la evaluación integral del desempeño del sistema, considerando indicadores de precisión, tiempo de respuesta, confiabilidad en el registro automático y estabilidad operativa. Los resultados obtenidos en el Capítulo~5 mostraron una precisión promedio del 85.2\,\%, un tiempo medio de detección de 0.32\,s por rostro, una tasa de registros automáticos correctos del 97.5\,\% y una disponibilidad del sistema del 99.0\,\%. Estos valores reflejan un desempeño sólido y equilibrado entre exactitud y eficiencia, permitiendo valorar de forma objetiva el comportamiento del sistema en condiciones reales de aula.

\textbf{Conclusión:} La evaluación de los resultados permitió confirmar que el sistema mantiene niveles de rendimiento coherentes con los parámetros definidos en la metodología, demostrando estabilidad y fiabilidad en el proceso de reconocimiento y registro de asistencias.

\section{Logros alcanzados con relación al Objetivo General}

El objetivo general de desarollar un sistema de reconocimiento facial para el control de asistencia de estudiantes de la Facultad Politécnica de la Universidad Nacional del Este fue alcanzado satisfactoriamente. Se desarrolló una solución funcional que automatiza el registro de asistencia de manera precisa, rápida y sin intervención manual.

El sistema emplea técnicas de visión por computadora y modelos de reconocimiento facial capaces de detectar y registrar rostros en tiempo real. Su funcionamiento demostró resultados consistentes en entornos de aula, reduciendo errores humanos y optimizando el proceso de control de asistencia.

La implementación comprobó la viabilidad del sistema propuesto, consolidándose como una herramienta moderna y adaptable para la gestión académica de la FPUNE.

\section{Análisis de la hipótesis}

La hipótesis de esta investigación establecía que:

\begin{quote}
“El desarrollo de un sistema de reconocimiento facial, basado en los modelos \textit{YOLOv8n-Face} y \textit{DeepFace} que permitirá registrar la asistencia de los estudiantes de la FPUNE, reduciendo los errores del control manual y optimizando el tiempo destinado al proceso de registro.”
\end{quote}

Los resultados obtenidos durante el desarrollo y la validación del sistema respaldan la veracidad de esta hipótesis. En las pruebas experimentales, el sistema alcanzó una precisión promedio del 85.2\,\%, con un tiempo medio de detección de 0.32\,s por rostro y una tasa de registros automáticos correctos del 97.5\,\%. Estos indicadores cuantitativos reflejan un desempeño estable y confiable en condiciones reales de aula, manteniendo una baja incidencia de falsos positivos (2.1\,\%) y una disponibilidad general del 99.0\,\%.

Desde una perspectiva estadística, los valores obtenidos se mantienen dentro de los parámetros establecidos en la metodología, lo que permite considerar al sistema como funcionalmente eficiente para la automatización del proceso de control de asistencia. Las variaciones observadas en las pruebas bajo condiciones de iluminación o ángulo no afectaron significativamente la precisión promedio, evidenciando la estabilidad del modelo ante factores ambientales.

En consecuencia, la hipótesis queda validada empíricamente, demostrando que la aplicación de técnicas de reconocimiento facial constituye una alternativa viable, eficiente y estadísticamente consistente para automatizar el registro de asistencia en la Facultad Politécnica de la Universidad Nacional del Este.


\section{Solución al problema de investigación}

El problema identificado en el Capítulo 1 señalaba que el control de asistencia en la Facultad Politécnica de la Universidad Nacional del Este se realizaba mediante métodos tradicionales, como el llamado de lista o el registro manual en papel, generando errores humanos, pérdida de documentos y dificultades en la organización de los datos.

La solución desarrollada —un sistema de asistencia basado en reconocimiento facial— responde directamente a esta problemática, permitiendo registrar de forma automática la asistencia de los estudiantes y almacenar la información de manera digital y ordenada.

La implementación del sistema:

\begin{itemize}
    \item Elimina la necesidad de registros manuales, reduciendo significativamente los errores humanos.  
    \item Acelera el proceso de control de asistencia, registrando la presencia de los estudiantes en segundos.  
    \item Facilita la gestión de datos mediante una interfaz que permite consultar, registrar y actualizar la información en tiempo real.  
    \item Ofrece un funcionamiento estable en condiciones comunes de aula, garantizando un registro confiable.  
\end{itemize}

En síntesis, el sistema no solo resuelve el problema técnico planteado, sino que también mejora la eficiencia y la precisión del proceso de control de asistencia, contribuyendo a una gestión académica más moderna y automatizada dentro de la FPUNE.


\section{Sugerencias para futuras investigaciones}

{\renewcommand{\labelitemi}{\textbullet}
\begin{itemize}
    \item \textbf{Integrar verificación por doble factor}, incorporando mecanismos adicionales como códigos QR, tarjetas inteligentes o tecnología NFC, con el fin de reforzar los procesos de autenticación en situaciones críticas. Esta ampliación permitiría validar la identidad del usuario a través de dos canales independientes, reduciendo la posibilidad de errores o suplantaciones. En contextos de evaluaciones, laboratorios o actividades de control estricto, esta verificación complementaria aumentaría la fiabilidad del sistema y contribuiría al cumplimiento de estándares de seguridad institucional.

    \item \textbf{Evaluar el desempeño del sistema con una base de datos más amplia y heterogénea}, que incorpore un número significativamente mayor de estudiantes, docentes y personal administrativo, abarcando variaciones demográficas como edad, género, etnia y diferentes condiciones lumínicas. Una ampliación de esta naturaleza permitiría realizar un análisis estadístico más robusto sobre la precisión y tolerancia del modelo, así como validar su generalización frente a escenarios reales más complejos. Además, abriría la posibilidad de aplicar métricas adicionales, como sensibilidad, especificidad o área bajo la curva ROC, para una evaluación más detallada del rendimiento global.

    \item \textbf{Adaptar el sistema a dispositivos móviles o cámaras IP} con transmisión remota y conectividad en red, permitiendo su funcionamiento en aulas híbridas, espacios abiertos o ambientes con conectividad limitada. Este tipo de implementación favorecería la escalabilidad y portabilidad del sistema, extendiendo su uso a sedes descentralizadas o actividades fuera del campus universitario. Asimismo, posibilitaría una integración con plataformas virtuales institucionales, facilitando el registro de asistencia en entornos de educación a distancia.

    \item \textbf{Aplicar técnicas de aprendizaje continuo o incremental}, de modo que el modelo de reconocimiento facial pueda actualizarse progresivamente con nuevas muestras sin requerir un reentrenamiento completo. Esta característica mejoraría la adaptabilidad del sistema ante cambios naturales en el rostro de los usuarios (como crecimiento de barba, uso de gafas o variaciones por edad), optimizando el uso de recursos computacionales. De igual forma, permitiría mantener actualizado el modelo a largo plazo, conservando un alto nivel de precisión sin necesidad de interrupciones operativas ni de grandes volúmenes de datos de entrenamiento.

    \item \textbf{Explorar el uso de técnicas de anonimización o protección de datos biométricos}, garantizando el cumplimiento de los principios de privacidad, ética y legislación vigente, tales como la Ley de Protección de Datos Personales. Futuras líneas de trabajo podrían orientarse a la implementación de mecanismos de cifrado y almacenamiento seguro de características faciales, con el objetivo de minimizar riesgos de exposición o uso indebido de información sensible. Asimismo, se sugiere el estudio de estrategias de anonimización mediante hash o encriptación homomórfica, de forma que el sistema pueda escalar institucionalmente sin comprometer la seguridad ni la confidencialidad de los datos recolectados.
\end{itemize}
}
